% Zdroj: PDF priklady-leto-2022-one.pdf (sešit S1, PGVVH+UI), fyzická str. 1. Značka: společné okruhy. Zařazeno: I2.
\section{Binární vyhledávací stromy}
\szzotazka{léto 2022}{1}{Binární vyhledávací stromy}

\noindent\textbf{Zadání.}\par
\begin{enumerate}
\item Definujte binární vyhledávací strom (BVS).
\item Implementujte vložení prvku do BVS (bez vyvažování) a~rozeberte časovou
  složitost v~nejlepším a~nejhorším případě.
\item Implementujte nalezení následníka, tj.\ pro daný klíč $x$ vrchol s~nejmenším
  klíčem ostře větším než $x$, a~rozeberte časovou složitost v~nejlepším
  a~nejhorším případě.
\end{enumerate}
Implementaci popište pseudokódem.

\medskip
\noindent\textbf{Náčrt řešení.}\par

\textbf{(1)} \emph{Binární vyhledávací strom} je binární strom, v~němž každý vrchol
nese klíč a~pro každý vrchol platí \emph{uspořádací vlastnost}: všechny klíče
v~levém podstromu jsou menší než klíč vrcholu a~všechny klíče v~pravém podstromu
jsou větší než klíč vrcholu.

\textbf{(2) Vložení.} Sestupujeme od kořene a~podle porovnání klíče jdeme doleva,
nebo doprava; nový prvek vložíme jako list na místo, kde sestup narazí na prázdný
podstrom.
\begin{lstlisting}
Insert(strom T, klic k):
  if T.koren == null:
    T.koren = new Uzel(k); return
  v = T.koren
  loop:
    if k == v.klic: return      // klic uz ve strome je
    if k < v.klic:
      if v.levy == null: v.levy = new Uzel(k); return
      else: v = v.levy
    else:                       // k > v.klic
      if v.pravy == null: v.pravy = new Uzel(k); return
      else: v = v.pravy
\end{lstlisting}
Vložení projde jednou cestou od kořene nejvýše k~listu, stojí tedy $O(h)$, kde $h$ je
hloubka stromu. Nejhorší případ: degenerovaný strom (liána, např.\ vzniklý vkládáním
setříděných klíčů) má $h=n-1$ a~nový klíč se připojí pod jeho dolní konec, tj.\ $\Theta(n)$.
Nejlepší případ: sestup skončí hned pod kořenem, tj.\ $\Theta(1)$ (např.\ do liány
$1,2,\dots,n$ s~kořenem $1$ vkládáme $0$). Uvažujeme-li nejlepší \emph{tvar} stromu
při nejhorší volbě klíče, je to vyvážený strom s~$h=\Theta(\log n)$, tedy $\Theta(\log n)$.

\textbf{(3) Následník.} Pokud má vrchol pravý podstrom, následník je jeho nejlevější
vrchol. Jinak stoupáme nahoru, dokud nejsme levým synem; tento rodič je následník
(neexistuje-li, $x$ byl maximum a~následník není).
\begin{lstlisting}
Successor(uzel x):
  if x.pravy != null:                 // nejmensi v pravem podstromu
    v = x.pravy
    while v.levy != null: v = v.levy
    return v
  v = x; r = x.rodic                  // vystup nahoru, dokud nejsem levy syn
  while r != null and v == r.pravy:
    v = r; r = r.rodic
  return r
\end{lstlisting}
Obě větve projdou jednu cestu stromem, takže složitost je $O(h)$. Nejhorší případ
$\Theta(n)$: ve stromu vzniklém vkládáním $1,n,n-1,\dots,2$ je následník kořene $1$
klíč $2$ na dně levé větve pod~$n$, cesta má $n-1$ hran. Nejlepší případ $\Theta(1)$:
pravý syn $x$ je list, nebo $x$ je levý syn bez pravého podstromu. Ve vyváženém stromu
je i~nejhorší případ jen $\Theta(\log n)$.

\medskip
\textit{(V~PDF bez náčrtu řešení; náčrt doplněn k~výkladu okruhu.)}
